\documentclass[oneside]{ctexart}
\setlength{\headheight}{14pt}
\usepackage{xcolor}
\usepackage{enumerate}
\usepackage{amsmath}
\usepackage{tikz-cd}
\usepackage{mathtools}
\usepackage{amssymb}
\usepackage{extarrows}
\usepackage{amsthm}

\newtheorem{definition}{定义}[section]
\newtheorem{example}{例}[section]
\newtheorem{theorem}{定理}[section] 
\newtheorem{axiom}{公理}[section]
\newtheorem{lemma}{引理}[section] 
\newtheorem{proposition}{命题}[section] 
\newtheorem{corollary}{推论}[section] 
\newtheorem{remark}{注}[section]

\DeclareMathOperator{\Tr}{Tr}
\DeclareMathOperator{\tr}{tr}
\DeclareMathOperator{\rank}{rank}
\DeclareMathOperator{\diag}{diag}
\DeclareMathOperator{\Ree}{Re}
\DeclareMathOperator{\Imm}{Im}
\DeclareMathOperator{\Ker}{Ker}

\usepackage{arydshln}
\usepackage{amsfonts} 
\usepackage{mathrsfs} 
\usepackage{bm}
\usepackage[T1]{fontenc}
\usepackage{fix-cm}
\usepackage{lmodern}
\usepackage{anyfontsize}
\usepackage{graphicx}
\usepackage[colorlinks=true, linkcolor=blue, urlcolor=blue, citecolor=blue]{hyperref}
\usepackage{geometry}
\geometry{
    a4paper,
    left=2.5cm,
    right=2.5cm,
    top=3cm,
    bottom=2cm
}

\title{神经网络的模型与应用作业论文阅读笔记\\[0.6ex]
\large  \textit{Representations of the First Hitting Time Density of an Ornstein-Uhlenbeck Process} }
\author{}
\date{}

\begin{document}

\maketitle

\section{论文主要内容与整体思路}

本文研究 Ornstein-Uhlenbeck（OU）过程第一次到达固定水平 $a$ 的时间
\[
    \sigma_a=\inf\{t>0:U_t=a\}
\]
的概率密度函数 $p_{x\to a}^{(\lambda)}(t)$.

文章主线很清楚:

\begin{enumerate}[(1)]
    \item 先定义 OU 过程与首次到达时;
    \item 再求首次到达时的拉普拉斯变换
          \[
              u_{a}^{\alpha}(x)=\mathbb{E}_x^{(\lambda)}[e^{-\alpha\sigma_a}] ;
          \]
    \item 然后把这个拉普拉斯变换分别反演成三种表示:
          \begin{enumerate}[(i)]
              \item 级数表示;
              \item 积分表示;
              \item Bessel bridge 表示;
          \end{enumerate}
    \item 最后讨论这三种表示的数值近似.
\end{enumerate}

\section{预备知识:主要定义与基本公式}

\subsection{OU过程的定义}

OU 过程满足随机微分方程
\begin{equation}\label{eq:ou}
    dU_t=dB_t-\lambda U_t\,dt,\qquad U_0=x\in\mathbb{R}.
\end{equation}

其中 $B_t$ 是标准布朗运动,$\lambda>0$.漂移项 $-\lambda U_t$ 表示过程有回到原点的趋势,因此 OU 过程是典型的均值回复过程.

\subsection{首次到达时}
对固定水平 $a$,定义首次到达时
\[
    \sigma_a=\inf\{t>0:U_t=a\}.
\]
如果它有密度,则记为
\[
    \mathbb{P}_x^{(\lambda)}(\sigma_a\in dt)=p_{x\to a}^{(\lambda)}(t)\,dt.
\]

\subsection{Doob变换}
由线性随机微分方程可解得
\[
    U_t=e^{-\lambda t}\left(x+\int_0^t e^{\lambda s}\,dB_s\right).
\]
再利用时间变换,可写成
\begin{equation}\label{eq:doob}
    U_t=e^{-\lambda t}\bigl(x+W_{\tau(t)}\bigr),\qquad
    \tau(t)=\frac{e^{2\lambda t}-1}{2\lambda}.
\end{equation}
这说明 OU 过程和布朗运动本质上联系很深,后面许多 hitting time 公式都依赖这个表示.

\subsection{生成元与边值问题}

OU 过程的生成元为
\[
    \mathcal{A}f(x)=\frac12 f''(x)-\lambda x f'(x).
\]

对拉普拉斯变换
\[
    u_{a}^{\alpha}(x)=\mathbb{E}_x^{(\lambda)}[e^{-\alpha\sigma_a}],
\]
它满足边值问题
\begin{equation}\label{eq:sl}
    \begin{cases}
        \mathcal{A}u(x)=\alpha\,u(x), & x<a, \\
        u(a)=1,                              \\
        \lim\limits_{x\to-\infty}u(x)=0.
    \end{cases}
\end{equation}

这一步非常关键,因为它把随机过程问题转化成了一个二阶微分方程边值问题.

\subsection{拉普拉斯变换的显式公式}

文章给出的核心公式是
\begin{equation}\label{eq:laplace}
    u_{a}^{\alpha}(x)
    =
    \frac{H_{-\alpha/\lambda}(-x\sqrt{\lambda})}
    {H_{-\alpha/\lambda}(-a\sqrt{\lambda})}
    =
    \frac{e^{\tfrac{\lambda x^2}{2}}D_{-\alpha/\lambda}(-x\sqrt{2\lambda})}
    {e^{\tfrac{\lambda a^2}{2}}D_{-\alpha/\lambda}(-a\sqrt{2\lambda})}.
\end{equation}

其中 $H_\nu$ 是 Hermite 函数,$D_\nu$ 是抛物柱函数.

文章还给出
\begin{equation}\label{eq:scaling}
    p_{x\to a}^{(\lambda)}(t)
    =
    \lambda\,p_{x\sqrt{\lambda}\to a\sqrt{\lambda}}^{(1)}(\lambda t).
\end{equation}

故只需研究 $\lambda=1$ 的情形.

\section{级数表示}

\subsection{思路}

这一部分的思想是:

\begin{enumerate}[(1)]
    \item 把拉普拉斯变换 $u_{a}^{\alpha}(x)$ 看成复变量 $\alpha$ 的函数;
    \item 研究它的极点;
    \item 再利用留数定理做反拉普拉斯变换.
\end{enumerate}

极点来自分母
\[
    D_{-\alpha/\lambda}(-a\sqrt{2\lambda})=0.
\]
设相应零点为 $\nu_1,\nu_2,\dots$,则密度可以展开成一串指数函数之和.

\subsection{主要公式}

级数表示为
\begin{equation}\label{eq:series}
    p_{x\to a}^{(\lambda)}(t)
    =
    -\lambda e^{\tfrac{\lambda(x^2-a^2)}{2}}
    \sum_{j=1}^{\infty}
    \frac{D_{\nu_{j,-a\sqrt{2\lambda}}}(-x\sqrt{2\lambda})}
    { D^{'}_{\nu_{j,-a\sqrt{2\lambda}}}(-a\sqrt{2\lambda})|_{\nu=\nu_j}}
    e^{-\lambda \nu_{j,-a\sqrt{2\lambda} }t}.
\end{equation}


\section{积分表示}

\subsection{思路}

这一部分是把拉普拉斯变换做解析延拓,再转成余弦反演.核心思想是:

\begin{enumerate}[(1)]
    \item 先把参数 $\alpha$ 推广到复数;
    \item 取纯虚数参数,得到 Fourier/余弦型表示;
    \item 再反演得到密度.
\end{enumerate}

\subsection{主要公式}

文章得到
\begin{equation}\label{eq:integral}
    p_{x\to a}^{(\lambda)}(t)
    =
    \frac{2\lambda}{\pi}
    \int_0^\infty
    \cos(\alpha\lambda t)
    \frac{
        Hr_{-\alpha}(-\sqrt{\lambda}a)Hr_{-\alpha}(-\sqrt{\lambda}x)
        +
        Hi_{-\alpha}(-\sqrt{\lambda}x)Hi_{-\alpha}(-\sqrt{\lambda}a)
    }{
        \left(Hr_{-\alpha}(-\sqrt{\lambda}a)\right)^2
        +
        \left(Hi_{-\alpha}(-\sqrt{\lambda}a)\right)^2
    }\,d\alpha.
\end{equation}

这里 $Hr_\alpha,Hi_\alpha$ 分别表示 Hermite 函数在纯虚参数下的实部和虚部.

\section{Bessel bridge表示}

\subsection{思路}

这一部分采用概率方法.核心思想是:

\begin{enumerate}[(1)]
    \item 先用 Girsanov 定理把 OU 过程和布朗运动联系起来;
    \item 再把布朗运动在首次到达条件下的路径,与三维 Bessel bridge 联系起来;
    \item 最后把密度写成桥过程某个泛函的期望.
\end{enumerate}

\subsection{布朗运动首次到达时密度}

当 $\lambda=0$ 时,有经典公式
\begin{equation}\label{eq:bm}
    p_{x\to a}^{(0)}(t)
    =
    \frac{|a-x|}{\sqrt{2\pi t^3}}
    \exp\left(-\frac{(a-x)^2}{2t}\right).
\end{equation}

\subsection{主要公式}

文章得到
\begin{equation}\label{eq:bessel}
    p_{x\to a}^{(\lambda)}(t)
    =
    e^{-\tfrac{\lambda(a^2-x^2-t)}{2}}
    \,p_{x\to a}^{(0)}(t)\,
    \mathbb{E}_{0\to a-x}
    \left[
        \exp\left(
        -\frac{\lambda^2}{2}\int_0^t (r_s-a)^2\,ds
        \right)
        \right].
\end{equation}

其中 $r_s$ 是从 $0$ 到 $a-x$ 的三维 Bessel bridge.


\section{数值近似的基本思路}

文章最后比较了三种数值方法:

\begin{enumerate}[(1)]
    \item \textbf{级数法}:截断 \eqref{eq:series} 的前若干项;
    \item \textbf{积分法}:对 \eqref{eq:integral} 做数值积分;
    \item \textbf{Monte Carlo 法}:对 \eqref{eq:bessel} 中的 Bessel bridge 泛函做模拟.
\end{enumerate}

总体上:

\begin{itemize}
    \item 级数法和积分法属于解析方法,精度较高;
    \item Bessel bridge 方法属于概率方法,更灵活,但数值误差通常更大.
\end{itemize}
\end{document}
